\section{Канон вполне-фундированности}\label{sec:wellfounded}

Двигатель --- один запрет: нет бесконечной строго убывающей цепи на вполне-фундированном носителе
(\code{no\_perpetual\_engine\_of\_wellFounded}, \cref{sec:engine}) --- и он переносится по \emph{любому}
рангу в вполне-фундированный порядок (\code{no\_perpetual\_engine\_of\_rank}). Каждый фронт до сих пор
был следствием этой единственной теоремы. Направленный на классический \emph{канон завершаемости}, чьё
всё содержание есть <<бесконечного спуска нет>>, он даёт зелёные инстансы, реюзая mathlib; ни один не
добавляет аксиомы, и таинт не двигается.

\begin{theorem}[\code{hydra\_no\_perpetual\_engine}]\label{thm:hydra_no_perpetual_engine}
\green{} Для любого вполне-фундированного $r$ отношение среза Кирби--Пэриса \code{Relation.CutExpand}~$r$
не несёт вечного двигателя: битва гидры завершается.
\end{theorem}

\emph{Идея доказательства.} Срез --- один шаг \code{CutExpand}; mathlib-евский \code{WellFounded.cutExpand}
даёт $\mathrm{WellFounded}(r) \Rightarrow \mathrm{WellFounded}(\code{CutExpand}\,r)$, что подаётся в
\code{no\_perpetual\_engine\_of\_wellFounded}.

\begin{theorem}[\code{omega\_worm\_no\_perpetual\_engine}]\label{thm:omega_worm_no_perpetual_engine}
\green{} Шаг на состоянии $(a,b)$ роняет большой регистр $a$ на единицу, а малый регистр $b$ может
прыгнуть на любое значение; процесс всё равно не несёт вечного двигателя --- под \emph{ординальным}
рангом $\omega\cdot a + b$.
\end{theorem}

\emph{Идея доказательства.} Хотя $b$ может взрываться, $\omega\cdot a + k < \omega\cdot(a+1) + b'$ при
любом $k$ (ибо $k < \omega$), так что \code{no\_perpetual\_engine\_of\_rank} с областью значений
\code{Ordinal} (а не $\mathbb{N}$) это закрывает --- первый трансфинитный ранг в программе и механизм за
теоремой Гудстейна и гидрой: значение взрывается, ординал падает. Полный слабый Гудстейн (цифры по
основанию $b$, прочитанные как ординал ниже $\omega^\omega$) продолжает тот же ранг; здесь не
формализован.

\begin{theorem}[\code{markov\_no\_perpetual\_engine}]\label{thm:markov_no_perpetual_engine}
\green{} Дерево Маркова $x^2+y^2+z^2=3xyz$ не несёт вечного двигателя прыжков Виеты $z \mapsto 3xy-z$:
высота $x+y+z$ строго убывает (\code{markov\_vieta\_lt}), так что дерево сводится к корню.
\end{theorem}

\emph{Идея доказательства.} При $x \le y \le z$ и $y < z$ прыжок посылает $z$ в $z' = 3xy-z \le y < z$
(ибо $y$ лежит между корнями $t^2 - 3xy\,t + x^2 + y^2$), так что верх и высота падают; замкнутость
Виеты (\code{isMarkov\_vieta}) держит тройку марковской. Это \emph{маска}: половина завершаемости зелена,
но гипотеза единственности Фробениуса (всякое число Маркова --- верх \emph{единственной} тройки, открыта
с 1913) спуска не имеет и не затронута --- зелёная стена, честные красные ворота.

\begin{theorem}[\code{poincare\_dividing\_line}]\label{thm:poincare_dividing_line}
\green{} На пространстве с мерой: всякое вполне-фундированное отношение запрещает вечный двигатель
(спуск, а значит и всякий возврат), тогда как консервативное --- например, сохраняющее конечную меру,
обратимое --- отображение возвращается по Пуанкаре, и почти всякая точка множества возвращается в него
бесконечно часто.
\end{theorem}

\emph{Идея доказательства.} Первая половина --- двигатель; вторая --- mathlib-евский
\code{Conservative.ae\_mem\_imp\_frequently\_image\_mem}. Вместе они --- те противоположности, что
\code{universal\_engine\_dividing\_line} называет на $\mathbb{R}$, теперь на динамической системе: строгий
спуск не возвращается, сохраняющая меру обратимость --- всегда --- замыкая ось обратимости квантового
прочтения (\cref{sec:coda}).

\begin{remark}[фундированность, а не самоссылка]
Этот канон вынуждает одну поправку честности. Теорема Гудстейна и гидра --- \emph{хрестоматийные}
примеры феномена гёделевой неполноты (истинны, но недоказуемы в арифметике Пеано), так что <<не
гёделев феномен>> надо читать точно: наш запрет --- принцип \emph{фундированности}, а не
\emph{самоссылки}. Гёделева фраза строится диагонализацией; эти же независимы потому, что их
доказательствам нужна трансфинитная индукция до $\varepsilon_0$, а <<$\varepsilon_0$ вполне-фундирован>>
--- ровно тот порядок, что PA доказать не может (Генцен: proof-theoretic ordinal PA есть
$\varepsilon_0$). Двигатель --- вполне-фундированная машина \emph{за} естественными независимостями,
достигаемая фундированностью, а не самоссылкой; <<не гёделев>> --- утверждение о \emph{механизме}, а не
бегство от неполноты (см.\ \cref{sec:coda}).
\end{remark}
